周报上写着：新模型准确率 \(0.912\)。

两周后复盘，这个数字变成了 \(0.897\)。模型退化了吗？会上没人能回答，因为大家不清楚这两个数字里有多少属于模型、多少属于这次抽到的样本。如果换一批同样规模的测试数据，\(0.912\) 本来就会变成 \(0.905\) 或 \(0.921\)——那么 \(0.897\) 到底算不算异常？

这就是数理统计要处理的方向。概率论从已知分布出发，问数据会怎样；统计反过来——手上只有一次有限样本，要说出关于生成机制的话。而这个方向天然是\textbf{欠定}的：同一份数据可以由许多不同的机制产生，所以任何结论都必须带着它的不确定性一起交付。

一句负责任的统计结论因此有三个部件：一个数、一个不确定范围、以及使这两者成立的一组假设。三个部件缺一不可，而实践中被省略的通常是第三个。

\begin{center}
\begin{tikzpicture}[x=1cm,y=1cm,font=\small,>=stealth]
  \draw[black!45,fill=blue!9] (0,1.5) rectangle (3.4,2.9);
  \node[align=center] at (1.7,2.2) {一次\\有限样本};
  \draw[->,black!55] (3.5,2.2) -- (4.6,2.2);
  \draw[black!45,fill=green!8] (4.7,3.35) rectangle (12.9,4.25);
  \node[align=left,anchor=west] at (4.9,3.8) {点估计：给出一个数（第二、三章）};
  \draw[black!45,fill=green!8] (4.7,2.2) rectangle (12.9,3.1);
  \node[align=left,anchor=west] at (4.9,2.65) {区间与检验：给出不确定范围（第四、五章）};
  \draw[black!45,fill=green!8] (4.7,1.05) rectangle (12.9,1.95);
  \node[align=left,anchor=west] at (4.9,1.5) {重采样：让数据自己评估整条流程（第六章）};
  \draw[black!45,fill=orange!10] (4.7,-0.1) rectangle (12.9,0.8);
  \node[align=left,anchor=west] at (4.9,0.35) {决策：把概率变成行动（第七章）};
  \draw[->,black!55] (4.6,2.25) -- (4.65,3.8);
  \draw[->,black!55] (4.6,2.1) -- (4.65,1.5);
  \draw[->,black!55] (4.6,2.0) -- (4.65,0.35);
  \node[align=center,font=\footnotesize,black!65] at (1.7,0.35) {每一步都挂着\\一组假设——\\它们是结论\\有效的条件};
\end{tikzpicture}

\emph{从一个数走到一个行动，中间每加一层承诺，就要多付一组假设。这一部分的大半篇幅在讲这些假设分别防住什么，以及它们塌掉时结论会以什么方式出错。}
\end{center}

\subsection{五条贯穿始终的判断}

\textbf{参数固定，统计量随机。}真实准确率是一个未知的常数，\(0.912\) 是它的一次实现。一切频率派的不确定性都来自这个随机性，而抽样分布——如果重复整个实验，这个数会怎样分布——就是把一次观测变成一句带范围的断言的桥梁。它是个思想实验：我们只有一份数据，它描述的是没有发生的那些平行世界。bootstrap 之所以有用，正是因为它用经验分布把这个思想实验重放了一遍。

\textbf{无偏不等于好。}均方误差拆成偏差平方加方差，而有偏但方差小的估计量可以整体更优。这解释了正则化为什么在统计上也划算，也解释了 Cramér–Rao 界为什么只约束无偏估计——有偏估计能突破它，那个缺口正是收缩与正则利用的空间。

\textbf{置信度属于程序，不属于这次的参数。}\(95\%\) 说的是造区间这套程序在重复使用下有 \(95\%\) 的时候会盖住真值；对手上这一个区间，真值要么在里面要么不在。想说``参数落在区间里的概率''，那是可信区间的语言，需要先给出先验。同理，P 值是 \(P(\text{数据}\given H_0)\) 而不是 \(P(H_0\given\text{数据})\)——方向一旦搞反，后面全错。

\textbf{检验是一条决策规则，两类错误必须同时看。}只控制第一类错误而不问功效，等于只承诺``不冤枉好人''却不承诺``抓得住坏人''。功效不足时实验做了等于没做，而且更糟：低功效下能通过显著性筛子的效应量被系统性夸大。多做几次检验、或看到显著就停，会把名义上的错误率悄悄推高——这与鞅论里的可选停时是同一个漏洞。

\textbf{评估的是整条流程，不是那个模型。}如果特征选择、标准化、超参搜索在划分之前做，信息就泄漏了，交叉验证会给出系统性乐观的估计。正确做法是把整条流程放进每一折内部重做；一个便宜的自检是把标签打乱重跑，指标应当回到随机水平。

\subsection{十章怎么安排}

\hyperref[ch:086]{``从一次数据到重复抽样''}是地基，它几乎不含计算，却决定了后面所有话该怎么说。\hyperref[ch:087]{``点估计：构造规则并评价其长期表现''}给出估计量的构造与评价，其中最大似然与 MAP 两篇直接对应训练目标与正则化。

\hyperref[ch:089]{``区间估计：让不确定性进入答案''}与\hyperref[ch:090]{``假设检验：在两类错误之间设计规则''}是使用频率最高的两章，也是误读最集中的两章。做实验分析的读者请把这两章读完整，尤其是 P 值那一篇。

\hyperref[ch:091]{``重采样与模型选择：让数据评估整个流程''}最贴近日常工程，bootstrap、置换检验、交叉验证三件工具加上信息泄漏这条警告，价值很高。\hyperref[ch:092]{``统计决策与可信预测：模型给出概率之后怎样行动''}接着回答部署阶段的问题——阈值怎么定、概率可不可信、容量有限时怎么排序，其中 conformal prediction 给出的是有限样本、无分布假设的覆盖保证。

\hyperref[ch:088]{``数据压缩与可估计信息''}偏理论但有实用回报：Fisher 信息把``参数辨认不出来''变成可诊断的量，Cramér–Rao 界则区分``换个估计量''与``这份数据本身就问不出更多''。

\hyperref[ch:093]{``回归与方差分析''}、\hyperref[ch:094]{``随机矩阵理论：高维协方差的谱失真''}、\hyperref[ch:095]{``GARCH 与波动率建模''}三章按需取用。随机矩阵那章值得一读：特征数与样本量可比时，样本协方差的谱会系统性失真，碎石图的肘部可能完全是噪声造出来的。

读这一部分时，把每个结论后面那句``在什么条件下''找出来。省略它不会立刻报错，只会让人对一个从未被验证过的前提充满信心。
